\documentclass[9pt]{beamer-control}
\usepackage{beamer-control-workshop}
\begin{document}
\CONCEPT[7]{Week 7: PID Control}

\begin{frame}
\frametitle{Introduction}
In this workshop, we will design and implement PID controllers for various control applications.
\end{frame}

\SUBCONCEPT{PID basics}

\begin{frame}[fragile]
\frametitle{PID in Matlab}
PID controllers can be defined in Matlab using the  $\texttt{pidstd(Kp,Ti,Td)}$ command
\includematlab{workshop7.m}{part1} 
Here we are setting $K_p=1$, $T_i=1$, and $T_d=1$. The PID object in Matlab is different to a transfer function, but the transfer function can be recovered using
\includematlab{workshop7.m}{part2} 
The following indexing syntax can also be used to view and change variable values
\includematlab{workshop7.m}{part3} 
\end{frame}


\begin{frame}[fragile]
	\frametitle{PID in Matlab}
	Using a transfer function we saw in last weeks workshop, let us try to control the closed loop system with a PID controller.
	\includematlab{workshop7.m}{part4} 
	Try to change the $K_p$, $T_i$, and $T_d$ variables from these initial values to achieve a better closed loop step response. \\[10pt]
	\alert{What does `better' mean in this context?}
	
\end{frame}

\begin{frame}[fragile]
	\frametitle{PID in Simulink}
	We can of course also design PID controllers in Simulink. In addition to blocks we have seen before, we will also need a Derivative block and a Transfer Function block. Try to recreate your Matlab step response in Simulink. \\
	Remember that 
	\[K_i = \frac{K_p}{T_i}, \quad K_d = K_pT_d\]
	when transferring PID gains from $\texttt{pidstd()}$ to Simulink or vice versa!\\[10pt]
	Your Simulink should look something like this (make sure you are incorporating negative feedback by appropriately changing the left-most Sum block!)
	\includegraphics[width=0.9\linewidth]{workshop7_simulink}
\end{frame}


\SUBCONCEPT{Pole placement}
\begin{frame}[fragile]
	\frametitle{Placing poles with PID}
	Let us follow an example from the notes for how we can place poles using PID.\\
	A generic first-order system is given by the transfer function 
	\[P(s) = \frac{b}{s+a}\]
	Comparing the desired and closed loop characteristic polynomials with a PI controller
	\[s^2+a_1s+a_2 = s^2 (a_bk_p)s+bk_i\]
	We may therefore assign arbitrary values of the closed loop poles via PI control ($K_d=0$)
	\[K_p=\frac{a_1-a}{b}, \quad K_i = \frac{a_2}{b}\]
\end{frame}


\begin{frame}[fragile]
	\frametitle{Placing poles with PID}
	Consider the system
	\[P(s) = \frac{3}{s-2}\]
	which is unstable in open loop.\\[10pt]
	
	Using the previous slide, derive a PI controller to place the closed loop poles at $\lambda = -1, -5$. \\[10pt]
	
	Test your working by constructing the closed loop system with your PID controller in Matlab, and determining the pole locations using  $\texttt{pole(sys)}$ \\[10pt]
	
	\alert{How robust do you think this process is for real systems?}
\end{frame}

\end{document}
